\chapter{Einleitung}

\section{Einführung}

\todo Samuel

\section{Team}

\begin{figure}[H]
  \centering
  \includegraphics[width=5cm]{images/securitysquad}
  \caption{Secutitysquad - Logo}
  \label{fig:securitysquad-logo}
\end{figure}

\section{Aufgabenstellung}
Anbieter von zwielichtigen Web-Angeboten greifen ihre User mit diversen Client-seitigen Methoden an. Beispiele für solche Angriffe sind Malware Downloads, Phishing, JavaScript Intranet Angriffe, oder Browser Exploits.

Ziel der Arbeit ist eine systematische Untersuchung der Aktivitäten von semi-legalen Webseiten im \ac{WWW}. Das erwartete Ergebnis ist ein Prüfportal, auf dem jene Webseiten automatisiert analysiert werden und Ergebnisse präsentiert werden sollen.

Nach dem ersten Schaffen einer Übersicht von interessanten Zielen, wie z.B. One-Click-Hoster oder File-sharing Sites sollen ausgewählte Webseiten manuell untersucht werden. Außerdem sollen verschiedene Angriffsszenarien zur weiteren Prüfung ausgewählt werden. Der Untersuchungsprozes der Webseiten soll im Verlauf dieser Arbeit stückweise automatisiert und in den Rahmen einer Prüfanwendung gebracht werden.

Abschließend sollen eine Vielzahl von Webseiten mit der Anwendung getestet und die Ergebnisse ausgewertet und dokumentiert werden.

\section{webifier}

\begin{figure}[H]
  \centering
  \includegraphics[width=5cm]{images/webifier}
  \caption{webifier - Logo}
  \label{fig:webifier-logo}
\end{figure}

webifier ist eine Anwendung, mit der Webseiten auf deren Seriosität und mögliche clientseitige Angriffe auf den Nutzer geprüft werden können. Sie besteht aus mehreren eigenständigen Teilanwendungen. Im Zentrum steht der Tester, welcher die einzelnen Tests verwaltet, ausführt und anschließend die Ergebnisse auswertet. Die Platform ist eine Webanwendung welche den Endnutzern eine grafische Oberfläche zur Verfügung stellt, um Webseiten zu überprüfen. Im Hintergrund setzt die Plattform auf den Tester auf.

\todo Samuel

\begin{itemize}
  \item Tests
  \item Data
  \item Statistics
\end{itemize}}

Um die Techniken und Algorithmen von webifier verstehen zu können sind einige Grundlagen erforderlich, welche nun im nächsten Kapitel genauer vorgsetellt werden.
